\subsection{Эмпирик үнэлгээний арга зүй}
\label{subsec:empirical_framework}

Энэхүү судалгаанд ковариацын матрицын үнэлэгчдийн багцын оновчлолд үзүүлэх нөлөөг түүврээс гадуурх (out-of-sample, OOS) гүйцэтгэлд тулгуурлан харьцуулан үнэлнэ. Үүнд түүврийн ковариацын матриц, тогтмол вариацын, диагональ, тогтмол корреляцын болон нэг хүчин зүйлийн зорилтот матрицад суурилсан шугаман агшилтын үнэлэгчид, мөн аналитик шугаман бус агшилтын үнэлэгчийг авч үзнэ.

Ковариацын матрицын үнэлгээний нарийвчлал болон түүнийг ашиглан байгуулсан багцын гүйцэтгэл нь ялгаатай ойлголтууд юм. Тухайлбал, ковариацын матрицын үнэлгээний алдааг бууруулсан арга нь багцын оновчлолын үр дүнд заавал хамгийн бага бодит эрсдэл бий болгох албагүй. Иймээс ковариацын үнэлэгчдийг зөвхөн үнэлсэн матрицын шинж чанараар бус, тэдгээрт үндэслэн байгуулсан багцын түүврээс гадуурх гүйцэтгэлээр харьцуулах шаардлагатай.

Ковариацын матрицын агшилтын үнэлэгчдийг багцын оновчлолд ашиглаж, түүврээс гадуурх эрсдэлийг харьцуулан үнэлэх аргачлалыг \citep{LedoitWolf2003} судалгаанд хэрэглэсэн. Мөн \citep{DeMiguelGarlappiUppal2009} судалгаанд багцын оновчлолын аргуудыг сарын өгөгдөлд суурилсан шилжих цонхны аргаар үнэлж, тэдгээрийн түүврээс гадуурх гүйцэтгэлийг харьцуулсан.

Эдгээр судалгааны эмпирик үнэлгээний зарчимд тулгуурлан энэхүү судалгаанд бүх үнэлэгчийн хувьд ижил хөрөнгө оруулалтын хязгаарлалттай long-only GMVP загварыг ашиглана. Үнэлгээний хугацаа, багцын жинг шинэчлэх давтамж болон гүйцэтгэлийн шалгуурыг бүх аргын хувьд ижил байлгана. Ингэснээр багцын гүйцэтгэлд үүсэх ялгааг ковариацын матрицын үнэлэгчийн сонголттой холбон харьцуулан үнэлэх боломж бүрдэнэ.


% ----------------------------------------------------------
% 2.6.1. Rolling-Window Out-of-Sample Evaluation
% ----------------------------------------------------------

\subsubsection{Шилжих цонхны (rolling-window) аргаар түүврээс гадуурх үнэлгээ хийх}
\label{subsubsec:rolling_window}

Ковариацын матрицын үнэлэгчдийн багцын оновчлолд үзүүлэх нөлөөг үнэлэхийн тулд шилжих цонхны (rolling-window) аргыг ашиглана. Энэ арга нь тогтмол урттай түүхэн өгөгдөлд үндэслэн загварын параметрүүдийг үнэлж, багцын оновчтой жинг тодорхойлсны дараа уг жинг сургалтын өгөгдөлд ороогүй дараагийн хугацааны өгөөжид хэрэглэдэг.

\citep{DeMiguelGarlappiUppal2009} судалгаанд багцын оновчлолын аргуудыг харьцуулахдаа шилжих цонхны аргаар параметрүүдийг дахин үнэлж, дараагийн сарын өгөөжөөр түүврээс гадуурх гүйцэтгэлийг тодорхойлсон.

Нийт сарын өгөөжийн ажиглалтын тоог $n$, хөрөнгийн тоог $N$, сургалтын цонхны уртыг $T$ гэж тэмдэглэе. $t$ сарын $N$ хөрөнгийн өгөөжийн векторыг
\[
\boldsymbol{R}_t
=
(R_{1,t},R_{2,t},\ldots,R_{N,t})^\top
\]
гэж тодорхойлно.

$t$ сарын эцэст ашиглах сургалтын өгөгдөл нь
\[
\mathcal{D}_t
=
\left\{
\boldsymbol{R}_{t-T+1},
\boldsymbol{R}_{t-T+2},
\ldots,
\boldsymbol{R}_t
\right\},
\qquad
t=T,T+1,\ldots,n-1
\]
байна.

Ковариацын матрицын үнэлгээний аргыг $k=0,1,\ldots,5$ гэж тэмдэглэе. Энд $k=0$ нь түүврийн ковариацын матриц, харин $k=1,\ldots,5$ нь агшилтын таван үнэлэгчийг илэрхийлнэ.

Сургалтын цонх бүрд $k$-р аргаар тооцсон ковариацын матрицыг
\[
\widehat{\boldsymbol{\Sigma}}_t^{(k)}
=
\widehat{\boldsymbol{\Sigma}}^{(k)}
(\mathcal{D}_t)
\]
гэж тэмдэглэнэ.

Нэг хүчин зүйлийн зорилтот матрицад суурилсан үнэлэгчийн хувьд $\mathcal{D}_t$-тэй ижил хугацааны зах зээлийн индексийн өгөөжийг нэмж ашиглана. Бусад үнэлэгчийн хувьд тухайн сургалтын цонхны хөрөнгүүдийн өгөөжид үндэслэн ковариацын матрицыг үнэлнэ.

Үнэлсэн ковариацын матрицад үндэслэн long-only GMVP-ийн оновчтой жинг дараах бодлогоор тодорхойлно:
\[
\begin{aligned}
	\widehat{\boldsymbol{w}}_t^{(k)}
	&=
	\underset{\boldsymbol{w}}{\arg\min}
	\quad
	\boldsymbol{w}^{\top}
	\widehat{\boldsymbol{\Sigma}}_t^{(k)}
	\boldsymbol{w},\\
	\text{subject to}\quad
	&\boldsymbol{1}^{\top}\boldsymbol{w}=1,\\
	&w_i\geq 0,
	\qquad i=1,\ldots,N.
\end{aligned}
\]

Энд $\widehat{\boldsymbol{w}}_t^{(k)}$ нь $t$ сарын эцэст тодорхойлж, дараагийн сард хэрэглэх багцын жингийн вектор юм.

Тухайн жинг ашиглан $t+1$ сард хэрэгжих багцын өгөөжийг
\[
r_{p,t+1}^{(k)}
=
\left(
\widehat{\boldsymbol{w}}_t^{(k)}
\right)^\top
\boldsymbol{R}_{t+1}
\]
гэж тооцно.

Энэ тооцоололд $\boldsymbol{R}_{t+1}$ нь ковариацын матриц болон багцын жинг тодорхойлоход ашиглагдахгүй. Иймээс ирээдүйн мэдээллийг сургалтын өгөгдөлд оруулахгүй байх нөхцөл хангагдана.

Дараагийн алхамд сургалтын цонхыг нэг сараар урагшлуулж, ковариацын матрицын бүх үнэлэгч, тэдгээрийн агшилтын эрчим болон багцын оновчтой жинг шинээр тооцно. Энэ үйлдлийг үнэлгээний хугацааны турш давтан гүйцэтгэнэ.

Энэхүү судалгааны үндсэн шинжилгээнд $N=10$ хөрөнгө, $n=172$ сарын өгөөж, $T=120$ сарын сургалтын цонхыг ашиглана.

Түүврээс гадуурх өгөөжийн ажиглалтын тоог
\[
M=n-T
\]
гэж тэмдэглэвэл
\[
M=172-120=52
\]
байна.

Иймээс үнэлэгч бүрийн хувьд
\[
\left\{
r_{p,T+1}^{(k)},
r_{p,T+2}^{(k)},
\ldots,
r_{p,n}^{(k)}
\right\}
\]
гэсэн $52$ сарын өгөөжийн цуваа үүснэ.

Ийнхүү бүх аргыг ижил хугацааны бодит өгөөжөөр үнэлэх нь тэдгээрийн гүйцэтгэлийг харьцуулах үндэслэл болно.


% ----------------------------------------------------------
% 2.6.2. Portfolio Performance Measures
% ----------------------------------------------------------

\subsubsection{Багцын гүйцэтгэлийн үзүүлэлтүүд}
\label{subsubsec:performance_measures}

Ковариацын матрицын үнэлэгчдийн гүйцэтгэлийг түүврээс гадуурх өгөөжийн хэлбэлзэл, түүврийн ковариацын матрицтай харьцуулсан вариацын бууралт болон Sharpe-ийн харьцаагаар үнэлнэ.

\paragraph{Түүврээс гадуурх хэлбэлзэл (Out-of-Sample Volatility).}

$k$-р үнэлэгчээр байгуулсан багцын түүврээс гадуурх дундаж сарын өгөөжийг
\[
\overline{r}_p^{(k)}
=
\frac{1}{M}
\sum_{t=T}^{n-1}
r_{p,t+1}^{(k)}
\]
гэж тодорхойлно.

Тэгвэл түүврээс гадуурх вариацын үнэлгээ нь
\[
\widehat{\sigma}_{p,k}^{\,2}
=
\frac{1}{M-1}
\sum_{t=T}^{n-1}
\left(
r_{p,t+1}^{(k)}
-
\overline{r}_p^{(k)}
\right)^2
\]
байна.

Үүнээс багцын сарын стандарт хазайлтыг
\[
\widehat{\sigma}_{p,k}
=
\sqrt{
	\widehat{\sigma}_{p,k}^{\,2}
}
\]
гэж тооцно.

Сарын өгөөжийн стандарт хазайлтыг жилийн хэмжээнд шилжүүлэхдээ
\[
\widehat{\sigma}_{p,k}^{\,ann}
=
\sqrt{12}\,
\widehat{\sigma}_{p,k}
\]
гэсэн томьёог ашиглана.

Энэхүү хөрвүүлэлт нь сарын өгөөжүүдийн хугацааны автоковариацыг тэг гэж үзсэн нөхцөлд үндэслэнэ.

Түүврээс гадуурх хэлбэлзэл бага байх нь тухайн ковариацын үнэлэгчээр байгуулсан багц харьцангуй бага эрсдэлтэй байгааг илэрхийлнэ.

\paragraph{Харьцангуй вариацын бууралт (Relative Variance Reduction).}

Агшилтын үнэлэгчээр байгуулсан багцын эрсдэлийг түүврийн ковариацын матрицад суурилсан GMVP-ийн эрсдэлтэй харьцуулахын тулд харьцангуй вариацын бууралтыг ашиглана.

Үүнийг
\[
RVR_k
=
1-
\frac{
	\widehat{\sigma}_{p,k}^{\,2}
}{
	\widehat{\sigma}_{p,0}^{\,2}
},
\qquad k=1,\ldots,5
\]
гэж тодорхойлно.

Энд $\widehat{\sigma}_{p,0}^{\,2}$ нь түүврийн ковариацын матрицад суурилсан GMVP-ийн түүврээс гадуурх вариац юм.

$RVR_k>0$ бол $k$-р агшилтын үнэлэгчээр байгуулсан багц суурь загвараас бага вариацтай, харин $RVR_k<0$ бол их вариацтай байгааг илэрхийлнэ.

Харьцангуй вариацын бууралтыг хувиар илэрхийлэхдээ
\[
RVR_k^{(\%)}
=
100\left(
1-
\frac{
	\widehat{\sigma}_{p,k}^{\,2}
}{
	\widehat{\sigma}_{p,0}^{\,2}
}
\right)
\]
гэж тооцно.

\paragraph{Sharpe-ийн харьцаа (Sharpe Ratio).}

Багцын эрсдэлд харьцуулсан өгөөжийг үнэлэхийн тулд Sharpe-ийн харьцааг ашиглана.

Ерөнхий тохиолдолд $t$ сарын эрсдэлгүй өгөөжийг $r_{f,t}$ гэж тэмдэглэвэл багцын илүүдэл өгөөж нь
\[
r_{e,t}^{(k)}
=
r_{p,t}^{(k)}-r_{f,t}
\]
байна.

Түүврээс гадуурх дундаж илүүдэл өгөөжийг
\[
\overline{r}_e^{(k)}
=
\frac{1}{M}
\sum_{t=T}^{n-1}
\left(
r_{p,t+1}^{(k)}
-
r_{f,t+1}
\right)
\]
гэж тодорхойлно.

Тэгвэл сарын Sharpe-ийн харьцаа нь
\[
\widehat{SR}_k
=
\frac{
	\overline{r}_e^{(k)}
}{
	\widehat{\sigma}_{e,k}
}
\]
байна.

Энд $\widehat{\sigma}_{e,k}$ нь түүврээс гадуурх сарын илүүдэл өгөөжийн стандарт хазайлт юм.

Энэхүү судалгааны тооцоололд сарын эрсдэлгүй өгөөжийг
\[
r_{f,t}=0
\]
гэж авч үзсэн.

Иймээс
\[
\overline{r}_e^{(k)}
=
\overline{r}_p^{(k)},
\qquad
\widehat{\sigma}_{e,k}
=
\widehat{\sigma}_{p,k}
\]
бөгөөд сарын Sharpe-ийн харьцааг
\[
\widehat{SR}_k
=
\frac{
	\overline{r}_p^{(k)}
}{
	\widehat{\sigma}_{p,k}
}
\]
гэж тооцно.

Сарын өгөөжүүдийн автоковариацыг тэг гэж үзсэн нөхцөлд жилийн Sharpe-ийн харьцааг
\[
\widehat{SR}_k^{\,ann}
=
\sqrt{12}\,
\widehat{SR}_k
\]
гэж хөрвүүлнэ.

Sharpe-ийн харьцаа өндөр байх нь нэгж эрсдэлд ногдох дундаж илүүдэл өгөөж харьцангуй өндөр байгааг илэрхийлнэ.

GMVP-ийн зорилго нь Sharpe-ийн харьцааг хамгийн их байлгах биш, харин багцын вариацыг хамгийн бага байлгах явдал гэдгийг анхаарах шаардлагатай. Иймээс Sharpe-ийн харьцааг багцын гүйцэтгэлийн нэмэлт үзүүлэлт болгон ашиглана.


% ----------------------------------------------------------
% 2.6.3. Portfolio Weight Stability and Turnover
% ----------------------------------------------------------

\subsubsection{Багцын жингийн тогтвортой байдал ба эргэлт}
\label{subsubsec:weight_stability}

Ковариацын матрицын үнэлэгчийн сонголт нь багцын эрсдэлээс гадна хөрөнгө оруулалтын жингийн бүтэц болон хугацааны явц дахь өөрчлөлтөд нөлөөлж болно.

Иймээс багцын жингийн бүтцийг төвлөрлийн үзүүлэлтээр, харин жингийн өөрчлөлтийг багцын эргэлтээр үнэлнэ.

\paragraph{Багцын эргэлт (Portfolio Turnover).}

Багцын эргэлт нь багцын жинг шинэчлэх үед худалдах болон худалдан авах шаардлагатай хөрөнгийн хэмжээг илэрхийлнэ.

$t$ сарын өгөөж хэрэгжихээс өмнө ашигласан багцын жинг $\widehat{\boldsymbol{w}}_{t-1}^{(k)}$ гэж тэмдэглэе.

$t$ сарын өгөөж хэрэгжсэний дараа, дахин тэнцвэржүүлэхээс өмнөх $i$-р хөрөнгийн жин нь
\[
\widetilde{w}_{i,t}^{(k)}
=
\frac{
	\widehat{w}_{i,t-1}^{(k)}
	(1+R_{i,t})
}{
	1+r_{p,t}^{(k)}
}
\]
байна.

Энд
\[
r_{p,t}^{(k)}
=
\sum_{i=1}^{N}
\widehat{w}_{i,t-1}^{(k)}
R_{i,t}
\]
нь тухайн сард хэрэгжсэн багцын өгөөж юм.

Ийнхүү хөрөнгүүдийн өгөөж ялгаатай байснаас шалтгаалан багцын жин анхны түвшнээс өөрчлөгдөнө.

$t$ сарын эцэст шинээр тодорхойлсон $\widehat{\boldsymbol{w}}_t^{(k)}$ жинд шилжихэд шаардагдах багцын эргэлтийг
\[
TO_t^{(k)}
=
\frac{1}{2}
\sum_{i=1}^{N}
\left|
\widehat{w}_{i,t}^{(k)}
-
\widetilde{w}_{i,t}^{(k)}
\right|
\]
гэж тодорхойлно.

Энд $1/2$ үржүүлэгч нь худалдах болон худалдан авах хэмжээг давхар тооцохоос сэргийлнэ.

Түүврээс гадуурх хугацааны дундаж эргэлтийг
\[
\overline{TO}^{(k)}
=
\frac{1}{M-1}
\sum_{t=T+1}^{n-1}
TO_t^{(k)}
\]
гэж тооцно.

Эхний түүврээс гадуурх сард хэрэглэсэн жинг дараагийн сарын жингийн өөрчлөлтийг тооцох эхлэл болгон ашиглана. Иймээс $M=52$ сарын OOS өгөөжид харгалзах багцын эргэлтийг $M-1=51$ удаагийн дахин тэнцвэржүүлэлтээр тооцно.

Багцын эргэлт бага байх нь хөрөнгө оруулалтын жинг шинэчлэх шаардлага харьцангуй бага байгааг илэрхийлнэ. Энэ үзүүлэлт нь арилжааны зардал үүсэх боломжийг үнэлэхэд ач холбогдолтой боловч энэхүү судалгаанд арилжааны бодит зардлыг багцын өгөөжөөс хасаагүй.

\paragraph{Багцын жингийн төвлөрөл (Weight Concentration).}

Багцын жингийн төвлөрлийг Herfindahl--Hirschman Index (HHI)-ээр хэмжинэ.

$t$ сарын эцэст $k$-р үнэлэгчээр байгуулсан багцын төвлөрлийг
\[
HHI_t^{(k)}
=
\sum_{i=1}^{N}
\left(
\widehat{w}_{i,t}^{(k)}
\right)^2
\]
гэж тодорхойлно.

Long-only GMVP-ийн хувьд
\[
\frac{1}{N}
\leq
HHI_t^{(k)}
\leq 1
\]
байна.

Бүх хөрөнгийн жин ижил, өөрөөр хэлбэл $w_i=1/N$ үед
\[
HHI_t^{(k)}
=
\frac{1}{N}
\]
болох бөгөөд багцын төвлөрөл хамгийн бага байна.

Харин бүх хөрөнгийг нэг хөрөнгөд байршуулсан үед
\[
HHI_t^{(k)}=1
\]
болно.

Түүврээс гадуурх хугацааны дундаж төвлөрлийг
\[
\overline{HHI}^{(k)}
=
\frac{1}{M}
\sum_{t=T}^{n-1}
HHI_t^{(k)}
\]
гэж тооцно.

$HHI$ өндөр байх нь багцын жин цөөн тооны хөрөнгөд төвлөрч байгааг илэрхийлнэ.

Багцын эргэлт болон төвлөрлийг хамтад нь авч үзсэнээр ковариацын үнэлэгчдийн багцын жингийн бүтэц болон хугацааны явц дахь өөрчлөлтийг харьцуулан үнэлнэ.


% ----------------------------------------------------------
% 2.6.4. Statistical Comparison
% ----------------------------------------------------------

\subsubsection{Ковариацын үнэлэгчдийн статистик харьцуулалт}
\label{subsubsec:statistical_comparison}

Түүврээс гадуурх гүйцэтгэлийн үзүүлэлтүүдийн ялгаа нь санамсаргүй хэлбэлзлээс үүссэн эсэхийг үнэлэхийн тулд статистик харьцуулалт хийнэ.

Энэхүү судалгаанд түүврийн ковариацын матрицад суурилсан GMVP-ийг суурь загвар болгон авч, агшилтын үнэлэгч тус бүрийн багцын гүйцэтгэлийг түүнтэй харьцуулна.

$k$-р агшилтын үнэлэгч болон түүврийн ковариацын матрицад суурилсан багцын тухайн сарын квадрат өгөөжийн ялгааг
\[
d_t^{(k)}
=
\left(r_{p,t}^{(k)}\right)^2
-
\left(r_{p,t}^{(0)}\right)^2,
\qquad
t=T+1,\ldots,n
\]
гэж тодорхойлно.

Энэ хэмжигдэхүүн нь хоёр багцын өгөөжийн тэг утгын орчим дахь квадрат хэмжигдэхүүний ялгааг илэрхийлнэ.

Түүний түүврээс гадуурх дундаж нь
\[
\overline{d}^{(k)}
=
\frac{1}{M}
\sum_{t=T+1}^{n}
d_t^{(k)}
\]
байна.

Дараах таамаглалуудыг шалгана:
\[
H_0:
E[d_t^{(k)}]=0,
\]
\[
H_1:
E[d_t^{(k)}]\neq 0.
\]

$H_0$ таамаглал нь хоёр багцын хүлээгдэж буй квадрат өгөөжийн ялгаа тэг байгааг илэрхийлнэ.

$\overline{d}^{(k)}<0$ бол $k$-р агшилтын үнэлэгчээр байгуулсан багцын дундаж квадрат өгөөж суурь загвараас бага байгааг илэрхийлнэ.

Хугацааны цувааны ажиглалтуудын хооронд автокорреляц болон гетероскедастик шинж илэрч болох тул дундаж ялгааны стандарт алдааг heteroskedasticity and autocorrelation consistent (HAC) аргаар үнэлнэ.

Үүний тулд
\[
d_t^{(k)}
=
\mu_d^{(k)}
+
\varepsilon_t^{(k)}
\]
гэсэн зөвхөн тогтмол гишүүнтэй регрессийн загварыг авч үзнэ.

Энд $\mu_d^{(k)}$ нь квадрат өгөөжийн ялгааны хүлээгдэж буй дундаж бөгөөд түүний үнэлгээ нь
\[
\widehat{\mu}_d^{(k)}
=
\overline{d}^{(k)}
\]
байна.

HAC стандарт алдааг тооцохдоо лагийн тоог
\[
L=3
\]
гэж сонгож, Newey--West төрлийн HAC ковариацын үнэлгээг ашиглана.

Тэгвэл шалгах статистик нь
\[
t_{HAC}^{(k)}
=
\frac{
	\overline{d}^{(k)}
}{
	\widehat{SE}_{HAC}
	\left(
	\overline{d}^{(k)}
	\right)
}
\]
байна.

Энд $\widehat{SE}_{HAC} (\overline{d}^{(k)})$ нь дундаж ялгааны HAC стандарт алдаа юм.

Статистик шалгалтыг хоёр талт хэлбэрээр гүйцэтгэж, HAC статистикийн асимптотик нормаль ойролцоололд үндэслэн $p$-утгыг тооцно.

Энэхүү шалгалтаар агшилтын үнэлэгчдийн квадрат өгөөжийн дундаж ялгаа статистикийн хувьд ач холбогдолтой эсэхийг үнэлнэ.

Гэхдээ квадрат өгөөжийн ялгаа нь төвлөрүүлсэн вариацын ялгаатай яг ижил хэмжигдэхүүн биш юм. Учир нь
\[
E[r_p^2]
=
\operatorname{Var}(r_p)
+
\left(E[r_p]\right)^2.
\]

Иймээс HAC шалгалтын үр дүнг түүврээс гадуурх стандарт хазайлт болон харьцангуй вариацын бууралтын үзүүлэлттэй хамтатган тайлбарлана.

Түүнчлэн суурь загвартай таван агшилтын үнэлэгчийг тус тусад нь харьцуулж байгаа тул олон удаагийн статистик шалгалтын асуудал үүсэж болохыг анхаарна.

Энэхүү судалгааны үндсэн HAC шалгалтаар гарсан $p$-утгууд нь олон харьцуулалтын засвар хийгдээгүй утгууд юм. Иймээс тэдгээрийг тайлбарлахдаа олон шалгалтын улмаас санамсаргүй ач холбогдолтой үр дүн илрэх боломжийг харгалзана.

Мөн түүврээс гадуурх хугацаа $M=52$ сар тул статистик шалгалтын чадал хязгаарлагдмал байж болохыг анхаарах шаардлагатай.


% ----------------------------------------------------------
% 2.6.5. Sensitivity to N/T
% ----------------------------------------------------------

\subsubsection{Хөрөнгийн тоо ба сургалтын ажиглалтын тооны харьцааны мэдрэмжийн шинжилгээ}
\label{subsubsec:nt_sensitivity}

Ковариацын матрицын үнэлгээний нарийвчлалд хөрөнгийн тоо болон сургалтын ажиглалтын тооны харьцаа чухал нөлөөтэй.

Энэхүү харьцааг
\[
c=\frac{N}{T}
\]
гэж тодорхойлно.

Энд $N$ нь хөрөнгийн тоо, $T$ нь ковариацын матрицыг үнэлэхэд ашигласан сургалтын ажиглалтын тоо юм.

$N/T$ харьцаа ихсэхэд нэг хөрөнгөд ногдох ажиглалтын тоо харьцангуй багасч, түүврийн ковариацын матрицын үнэлгээний алдаа нэмэгдэх боломжтой.

Харин $N/T$ харьцаа бага үед түүврийн ковариацын матриц илүү олон ажиглалтад тулгуурлах боловч үнэлгээний алдаа бүрэн арилна гэсэн үг биш юм.

Энэхүү судалгаанд хөрөнгийн тоог $N=10$ гэж тогтмол байлгаж, сургалтын цонхны уртыг өөрчлөх замаар $N/T$ харьцааны мэдрэмжийг шинжилнэ.

Үүнд
\[
T\in\{60,90,120\}
\]
гэсэн сарын сургалтын цонхуудыг авч үзнэ.

Харгалзах харьцаанууд нь
\[
\begin{aligned}
	\frac{N}{T}
	&=\frac{10}{60}
	\approx 0.1667,\\[4pt]
	\frac{N}{T}
	&=\frac{10}{90}
	\approx 0.1111,\\[4pt]
	\frac{N}{T}
	&=\frac{10}{120}
	\approx 0.0833
\end{aligned}
\]
байна.

Сургалтын цонхны урт ялгаатай үед түүврээс гадуурх үнэлгээний эхлэх хугацаа мөн ялгаатай болох тул гүйцэтгэлийг шууд харьцуулах нь хугацааны нөлөөг агуулж болно.

Иймээс мэдрэмжийн шинжилгээнд бүх сургалтын цонхны хувьд ижил түүврээс гадуурх хугацааг ашиглана.

Тодруулбал, хамгийн урт сургалтын цонх $T_{\max}=120$ үед үүсэх
\[
M=n-T_{\max}
\]
ажиглалттай түүврээс гадуурх хугацааг бүх хувилбарт ижил хэрэглэнэ.

Энэхүү судалгаанд
\[
M=172-120=52
\]
учраас сургалтын цонхны гурван хувилбарыг ижил $52$ сарын OOS хугацаанд харьцуулна.

Ингэснээр $T=60$, $T=90$ болон $T=120$ үед ковариацын матрицыг үнэлэхэд ашигласан түүхэн ажиглалтын тоо ялгаатай боловч багцын гүйцэтгэлийг нэг ижил зах зээлийн хугацаанд харьцуулна.

Сургалтын цонхны хувилбар бүрд зургаан ковариацын үнэлэгчийг дахин тооцож, long-only GMVP-ийн оновчтой жинг тодорхойлно.

Дараа нь түүврээс гадуурх хэлбэлзэл, харьцангуй вариацын бууралт, Sharpe-ийн харьцаа, багцын эргэлт болон жингийн төвлөрлийг харьцуулна.

Энэ шинжилгээнд хөрөнгийн тоо тогтмол байгаа учраас $N/T$ харьцааны өөрчлөлт нь зөвхөн сургалтын цонхны уртын өөрчлөлтөөс үүднэ.

Иймээс шинжилгээний үр дүнг $N/T$ харьцааны нөлөөг сургалтын цонхны уртын нөлөөнөөс бүрэн тусгаарлан үнэлсэн гэж тайлбарлахгүй.

Харин сургалтын ажиглалтын тоо өөрчлөгдөхөд ковариацын матрицын агшилтын аргуудын түүврээс гадуурх гүйцэтгэл хэрхэн өөрчлөгдөж байгааг эмпирик түвшинд үнэлнэ.